% Part: history
% Chapter: biographies
% Section: georg-cantor

\documentclass[../../../include/open-logic-section]{subfiles}

\begin{document}

\olfileid{his}{bio}{can}

\olsection{غيورغ كانتور}

\olphoto{cantor-georg}{Georg Cantor}

زعمت سيرة مبكرة لغيورغ كانتور (\textsc{gay}-org
\textsc{kahn}-tor) أنه وُلد وعُثر عليه على متن سفينة كانت مبحرة إلى
سانت بطرسبرغ في روسيا، وأن والديه مجهولان. غير أن هذا غير صحيح، وإن كان
قد وُلد بالفعل في سانت بطرسبرغ عام 1845.

نال كانتور الدكتوراه في الرياضيات من جامعة برلين عام 1867. وهو معروف
بعمله في نظرية المجموعات، ويُنسب إليه تأسيس نظرية المجموعات بوصفها مبحثاً
بحثياً مستقلاً. وكان أول من برهن وجود مجموعات لانهائية مختلفة الأحجام.
وقد أثارت نظرياته، ولا سيما نظريته في اللانهايات، جدلاً واسعاً بين
رياضيي عصره، وكان عمله موضع خلاف.

ارتبطت معتقدات كانتور الدينية بعمله الرياضي ارتباطاً لا ينفصم؛ حتى إنه
زعم أن الله أبلغه مباشرةً نظرية الأعداد المتناهية التجاوز. وعانى كانتور
في أواخر حياته مرضاً نفسياً. فابتداءً من عام 1894، وبوتيرة أشد في سنواته
الأخيرة، كان يُودع المستشفى. وأفضى النقد الشديد لعمله، ومنه خلافه مع
الرياضي ليوبولد كرونيكر، إلى الاكتئاب وفتور اهتمامه بالرياضيات. وفي نوبات
الاكتئاب كان كانتور ينصرف إلى الفلسفة والأدب، بل نشر نظريةً تقول إن
فرانسيس بيكون هو مؤلف مسرحيات شكسبير.

توفي كانتور في 6 يناير 1918 في مصح بمدينة هاله.

\begin{reading}
للاطلاع على سيرتين كاملتين لكانتور، انظر \citet{Dauben1990} و
\citet{Grattan-Guinness1971}. وتُعرض آراء كانتور الجذرية أيضاً في برنامج
إذاعة بي بي سي 4 \emph{تاريخ موجز للرياضيات} \citep{Sautoy2014}.
ولمن يود سماع نظريات كانتور في قالب راب، انظر \citet{Rose2012}.
\end{reading}

\end{document}
